\section{Coding 的爆发：从写代码到长程代理}

访谈中最密集的技术讨论之一，是 coding 与 agentic workflow 的爆发。姚顺宇认为，OpenClaw 一类产品在技术上并不令人惊讶，因为相关能力在更早的 Claude/Opus 阶段已经可以展示；它真正的价值是让更多人意识到：可以让模型控制多个工具、调用多个模型、聚合结果，并完成很长的 long-horizon 任务。

\begin{figure}[H]
\centering
\includegraphics[width=0.88\textwidth]{frame-coding-explosion.jpg}
\caption{Coding 爆发章节：讨论 OpenClaw、Manus、agentic coding 与模型能力自然外溢。\protect\footnotemark}
\end{figure}
\footnotetext{视频画面时间区间：00:42:57--00:43:25。该帧对应 coding 产品形态讨论的中段。}

\begin{knowledgebox}{读图：coding 爆发的证据类型}
本节图像用于时间定位，不展示 OpenClaw 或 Manus 的界面。有效证据来自访谈中的因果链：模型工具使用能力已经具备、产品把可能性展示出来、长程任务共识被触发。不要把人物帧误读为产品截图。
\end{knowledgebox}


\begin{importantbox}{从代码补全到 long-horizon agent}
早期 coding 模型主要生成一段代码；agentic coding 则要求模型理解代码库、规划多步修改、运行测试、解释错误、回滚或修复，并在长时间上下文中保持目标一致。能力差异不只在“会不会写函数”，而在工具调用、状态管理、错误恢复和任务收敛。
\end{importantbox}

\subsection{Manus、OpenClaw 与 wrapper 的壁垒问题}

姚顺宇没有把 Manus 与 OpenClaw 的差异解释成清晰的技术代际。他更关注一个商业问题：如果 wrapper 没有自己的模型壁垒，长期生存会很难。当前很多壁垒仍在模型侧；产品侧是否能形成数据飞轮，还没有被证明。除了 agentic coding，他认为真正 AI-native 且已成功的场景并不多。

\begin{knowledgebox}{什么是 AI-native 场景}
AI-native 场景不是“把 AI 加到旧产品里”，而是任务形态本身因为模型能力而改变。Chatbot 可看作搜索的交互式扩展；agentic coding 则更接近新形态，因为它把写代码、运行环境、调试、测试和项目管理连成一条工作流。
\end{knowledgebox}

\subsection{产品经理为什么暂时难被替代}

访谈中对 Claude Code、Cowork 等产品的讨论，指向一个重要判断：AI 产品经理的价值不再只是按钮摆放和 feature 排序，而是理解“人如何与 AI 协作”。好的产品形态会改变交互方式，就像短视频改变内容消费，Claude Code 改变软件工程入口。

\begin{warningbox}{不要把 wrapper 一概看低}
“wrapper 没有壁垒”是市场常见说法，但并不意味着所有 wrapper 都无价值。真正的问题是：它能否形成工作流锁定、数据反馈、用户习惯、团队执行和模型侧深度集成。没有这些，wrapper 容易被模型公司吸收；有这些，就可能变成新入口。
\end{warningbox}

\subsection{本章小结}

Coding 的爆发不是单点模型能力提升，而是模型能力、工具链、评估、环境和产品形态共同成熟的结果。OpenClaw/Manus 的讨论说明：产品展示可能比技术突破更早触发大众共识，但长期壁垒仍取决于模型、数据和工作流闭环。

